\documentclass[review]{elsarticle}
\usepackage[utf8]{inputenc}
\usepackage[T1]{fontenc}
\usepackage{amsmath,amssymb}
\usepackage{tabularx,booktabs}
\usepackage{url}
\usepackage{textcomp}
\usepackage{graphicx}
\usepackage{lineno}
\setlength{\emergencystretch}{2em}
\renewcommand{\thesection}{\Roman{section}}
\renewcommand{\thesubsection}{\thesection.\Alph{subsection}}
\renewcommand{\thetable}{\Roman{table}}
\renewcommand{\thesection}{\Roman{section}}
\renewcommand{\thesubsection}{\thesection.\Alph{subsection}}
\renewcommand{\thetable}{\Roman{table}}
\journal{Solar Energy}
\begin{document}
\begin{frontmatter}

\title{Oxidation-Budgeted Design of Lead-Free RbGeI$_{3}$ Solar Cells: Back-Surface-Field Enhancement, Realistic Loss Budgeting and SHAP Sensitivity From SCAPS-1D Optimization to Manufacturing Tolerances, Benchmarked Against CsGeI$_{3}$}
\begin{abstract}
The rapid rise of lead-halide perovskite solar cells has been tempered by lead toxicity, motivating a search for lead-free absorbers with comparable optoelectronics. Rubidium germanium iodide (RbGeI$_{3}$) offers a direct, tunable bandgap, strong absorption, and improved environmental compatibility. Here we report a SCAPS-1D optimization of an RbGeI$_{3}$ cell organized around its central liability, Ge$^{2+}$ oxidation. After screening eight ETL--HTL combinations, the FTO/\allowbreak{}TiO$_{2}$/RbGeI$_{3}$/CuI/\allowbreak{}Au architecture was carried forward on stability and cost grounds (a D2/NiO cross-check through the same pipeline reaches 28.16\%, so this stack is not the efficiency optimum). Eleven parameters were varied sequentially on one frozen deck; the consolidated device reaches 27.06\% PCE. A Cu$_{2}$O-like back-surface field raises the ideal point to 28.01\% and recovers about one point on the identical full-realism stack (20.78\% $\rightarrow$ 21.76\%). Mapping V$_{\mathrm{Ge}}$ vacancy thermodynamics onto the defect ladder converts air exposure into a fab spec ($E^f\gtrsim0.46$ eV holds the headline at 300 K equilibrium, 0.57--0.65 eV for anneal freeze-in; ambient falls below 20.76\%), with a 15.53\% current-state rung marking what first fabrication must beat. The solver pipeline is calibration-checked against experimental $\beta$-CsPbI$_{3}$ \cite{rWang}, and SHAP analysis recovers bulk defects, bandgap and electron mobility as the dominant levers.
\end{abstract}
\begin{keyword}
lead-free perovskite; RbGeI$_{3}$; SCAPS-1D; back surface field; oxidation budget; SHAP sensitivity analysis
\end{keyword}
\end{frontmatter}


\section{Introduction}

\subsection{Background and motivation}
Global electricity demand keeps climbing while fossil fuels still supply its bulk with a heavy carbon cost \cite{r1,r2}. Solar is the fastest-expanding renewable technology, with installed capacity surpassing 2.2 TW in 2024 and levelized costs still falling \cite{r3}; continued growth needs absorbers combining high efficiency with low environmental impact and cost \cite{r4}. Devices are commonly classed into three generations --- crystalline silicon, thin films, and emerging devices --- and halide perovskites lead the third generation with certified single-junction efficiencies above 26\% and tandems past 33\% \cite{r3}.

\subsection{Perovskite solar cells and the lead problem}
Halide perovskites (ABX$_{3}$) attract photovoltaics with strong absorption, tunable gaps, ambipolar transport with long diffusion lengths, and weak exciton binding \cite{r7,r8}, rising from sensitized and meso-superstructured architectures \cite{r9,r10}. The champion absorbers MAPbI$_{3}$ and FAPbI$_{3}$ hold certified records \cite{r3} but contain lead, whose life-cycle contamination motivates Sn, Ge, Bi and Sb substitutes \cite{r11,r12}. Tin drew first attention since Sn$^{2+}$ shares the Pb$^{2+}$ lone-pair configuration \cite{r11}, but Sn$^{2+}$ oxidizes readily to Sn$^{4+}$, creating deep defects and fast degradation \cite{r11,r12}. Germanium offers the same lone-pair structure \cite{r13}; among Ge absorbers RbGeI$_{3}$ stands out for a near-optimal direct gap ($\approx$1.3--1.4 eV), MAPbI$_{3}$-comparable absorption by first principles \cite{r13,r17,r18}, and good thermal stability \cite{r14}.

\subsection{Role of numerical simulation and research gap}
Because experimental optimization of multilayer devices is slow and expensive, numerical simulation screens materials before fabrication. SCAPS-1D (Ghent University) self-consistently solves Poisson, continuity and drift-diffusion equations \cite{r19,r20} and is the workhorse for lead-free design, as a recent review of 54 SCAPS lead-free studies confirms \cite{rReview} (see their Fig.~1).

The RbGeI$_{3}$ simulation baseline is Pindolia et al.\ at 9.77\% for FTO/\allowbreak{}TiO$_{2}$/RbGeI$_{3}$/NiO/\allowbreak{}Ag \cite{r14}; predictions since reach 27.21\% \cite{r22} and 26.47\% \cite{r23}, each optimizing only a subset of parameters. The TiO$_{2}$/CuI combination used here has received comparatively little attention despite CuI's documented advantages as an HTL.

\subsection{Contributions and paper organization}
We close this gap with a SCAPS-1D study of an FTO/\allowbreak{}TiO$_{2}$/RbGeI$_{3}$/CuI/\allowbreak{}Au cell built around one idea: Ge$^{2+}$ oxidation converted into a quantitative fab spec. It contributes: (i) an oxidation defect-tolerance map from V$_{\mathrm{Ge}}$ thermodynamics on a frozen-deck eleven-step optimization (27.06\% ideal, 20.78\% stacked window, D2 like-for-like cross-check at 28.16\%); (ii) a Cu$_{2}$O BSF variant whose trap-immunity mechanism is isolated by controlled interface attribution (+0.98 pp on the identical realistic stack); and (iii) held-out SHAP design rules plus experimental grounding (CsPbI$_{3}$ calibration refit, frozen database snapshot as context). Band lineups and primary-source-verified benchmarks support each claim.


\section{Literature Review and Research Gap}
SCAPS-1D, often combined with DFT, is the workhorse for lead-free design before fabrication. Supplementary Table S1 surveys recent Ge- and Sn-based studies; all values there were cross-checked against primary sources, adopting the J--V-implied efficiency where inconsistent. CuI is a wide-gap ($E_{\mathrm{g}}$ $\approx$ 3.1 eV) p-type conductor whose valence maximum aligns with RbGeI$_{3}$ for hole extraction \cite{r25,r26}, with single-crystal mobilities above 100 cm$^{2}$ V$^{-1}$ s$^{-1}$ and low-temperature solution processing \cite{r27}; yet few SCAPS studies pair it with RbGeI$_{3}$, and none runs an eleven-parameter optimization on this combination. Comparator structures span TiO$_{2}$, PCBM, C$_{60}$, NiO, CuI, CBTS, Spiro-OMeTAD and Ag/Au contacts; note the Loumachi ETL is C$_{60}$ (not ``C6'') with an Ag contact \cite{r22}.

\section{Device Structure and Material Selection}

\subsection{Proposed device architecture}
We use a conventional n-i-p planar heterojunction, FTO/\allowbreak{}TiO$_{2}$/RbGeI$_{3}$/CuI/\allowbreak{}Au: FTO front electrode, TiO$_{2}$ ETL, RbGeI$_{3}$ absorber, CuI HTL, Au back contact (Figure 1). Material choices are justified below; initial parameters are in Supplementary Table S2.

\begin{figure}[htbp]
\centering
\includegraphics[width=0.85\linewidth]{figs/fig2.png}
\caption{Schematic diagram of the proposed FTO/\allowbreak{}TiO$_{2}$/RbGeI$_{3}$/CuI/\allowbreak{}Au perovskite solar cell structure, showing the five functional layers and the direction of photogenerated carrier transport under AM 1.5G illumination.}
\end{figure}

\subsection{Material selection rationale}
Fluorine-doped tin oxide (FTO, 500 nm) is the front electrode for its transparency, conductivity, and thermal/chemical stability \cite{r28}, the standard low-cost TCO choice for this stack.
TiO$_{2}$ is the ETL: its 3.2 eV gap passes visible light, its conduction-band minimum gives downhill electron transfer, and its deep valence band blocks holes \cite{r30,r31}, forming stable perovskite interfaces despite lower mobility than ZnO/SnO$_{2}$ \cite{r32}.
RbGeI$_{3}$ is the absorber (direct 1.3--1.4 eV gap, MAPbI$_{3}$-comparable absorption \cite{r17,r18}, good thermal stability); screening used 400 nm and 1.4 eV from literature \cite{r14}, both subsequently optimized (Section V).
CuI is the HTL (conductive, transparent, 0.10 eV valence offset to RbGeI$_{3}$, low cost \cite{r25,r26,r33}); screening used 100 nm. Au is the back contact, ohmic to CuI \cite{r34}. 

\subsection{Initial material parameters}
Initial parameters from published experimental and numerical lead-free studies \cite{r14,r15,r16,r19,r30} are in Supplementary Table S2 (interfaces: Table S3). All eight configurations were screened at 400 nm absorber (Section V.A, Table II); the consolidated device uses the optimized 700 nm (Section V.G).





\section{Methodology}

\subsection{SCAPS-1D simulation framework}
The stack was simulated with SCAPS-1D v3.3.10 \cite{r19,r20}, solving Poisson plus electron/hole continuity self-consistently under drift-diffusion transport (see original references for details).
\begin{equation}\label{eq:poisson}
\frac{d}{dx}\left(\epsilon\frac{d\psi}{dx}\right) =
-q\left(p-n+N_{D}-N_{A}+\rho_{t}\right)
\end{equation}
\begin{equation}\label{eq:electron}
\frac{dJ_{n}}{dx} = q\,(R-G)
\end{equation}
\begin{equation}\label{eq:hole}
\frac{dJ_{p}}{dx} = q\,(G-R)
\end{equation}
where $\epsilon$ is the permittivity, $\psi$ the potential, n/p the carrier concentrations, $N_{\mathrm{D}}$/$N_{\mathrm{A}}$ the doping, $\rho_{\mathrm{t}}$ the trapped charge, $J_{n}$/$J_{p}$ the current densities, and G/R the generation/recombination rates. SCAPS-1D computes the J--V characteristics from the layer properties below.

\subsection{Simulation conditions}
All simulations were performed under the AM 1.5G spectrum at an incident power density of 100 mW/cm$^{2}$ (one sun, STC) with the operating temperature fixed at 300 K, using the initial material parameters of Supplementary Table S2 and the interface defect densities of Supplementary Table S3 unless otherwise stated. During each optimization step, only one parameter was varied while all remaining parameters were held constant; the general simulation conditions are summarized in Supplementary Table S4. Absorption uses the SCAPS square-root law $\alpha$ = A$\sqrt{h\nu-E_{\mathrm{g}}}$ with A = 10$^{5}$ cm$^{-1}$ at the electrical bandgap, with radiative recombination at the literature perovskite coefficient ($K_{\mathrm{rad}}$ $10^{-10}$ cm$^{3}$/s) included in all headline runs. Contacts are ideal ($R_{\mathrm{s}}$ = 0, $R_{\mathrm{sh}}$ = $\infty$) unless stated, so fill factors are upper bounds. All eleven sweeps ran sequentially on one frozen deck generation with carry-forward (Section~V.B); failed runs were retried and the survivors surface as \texttt{FAILED} rather than stale data. Scripting limits (front-reflection handling, C--V capture, resistance post-corrections) and the full reproducibility protocol are documented in Supplementary Section~S7.



\subsection{Performance parameters}
Performance was evaluated using the four key parameters extracted from the simulated J--V curve: open-circuit voltage (V$_{\mathrm{OC}}$), short-circuit current density (J$_{\mathrm{SC}}$), fill factor (FF), and power conversion efficiency (PCE), computed as
\begin{equation}\label{eq:pce}
\mathrm{PCE} = (V_{\mathrm{OC}} \times J_{\mathrm{SC}} \times \mathrm{FF})/P_{\mathrm{in}}
\end{equation}
with $P_{\mathrm{in}}$ = 100 mW/cm$^{2}$ under AM 1.5G.

\subsection{Sequential optimization procedure}
Starting from the initial device, one parameter at a time (OPAT) was varied while all others were held constant, as detailed in Section V.
The procedure comprises eleven optimization steps: (1) absorber thickness (0.3--1.0 $\mu$m); (2) absorber bulk defect density (10$^{12}$--10$^{18}$ cm$^{-3}$); (3) absorber dielectric constant (10--35, extended until the trend was mapped); (4) absorber electron affinity (3.7--4.2 eV); (5) TiO$_{2}$/RbGeI$_{3}$ interfacial defect density (10$^{12}$--10$^{18}$ cm$^{-2}$); (6) RbGeI$_{3}$/CuI interfacial defect density (10$^{12}$--10$^{18}$ cm$^{-2}$); (7) absorber bandgap (1.2--1.7 eV, with 1.2 probed); (8) TiO$_{2}$ ETL thickness (0.01--0.09 $\mu$m); (9) CuI HTL thickness (0.05--1.9 $\mu$m); (10) conduction-band effective density of states N$_{\mathrm{C}}$ (1$\times$10$^{18}$--1$\times$10$^{20}$ cm$^{-3}$, literature floor); (11) valence-band effective density of states N$_{\mathrm{V}}$ (same range). Each step's value was carried forward under a triple-gate adoption rule: sweep maximum, literature-parameter compliance, and Shockley--Queisser admissibility with margin. Parameters that cannot be ordered from a supplier (dielectric constant) or tuned below physical floors (defect densities, N$_{\mathrm{C}}$/N$_{\mathrm{V}}$) were pinned at literature values with the full trace disclosed, rather than carried at grid edges; the consolidated set is summarized in Table IV.
Because each per-step value is identified with all other parameters fixed at prior-step values, the consolidated set may not be a global optimum. To test for missed cross-interactions, Section V.G reports a genetic-algorithm joint search over the six most sensitive parameters, which returns a best-found joint point (not a proven global maximum).

\subsection{Device configurations screened}
Eight planar configurations combining two ETLs (TiO$_{2}$, PCBM) with four HTLs (NiO, CuI, CBTS, Spiro-OMeTAD) around the RbGeI$_{3}$ absorber (D1--D8, Table II) were simulated at the initial absorber thickness of 400 nm (Supplementary Table S2) to permit a fair architecture comparison before the thickness optimization of Section V.B.

\begin{table}[htbp]
\centering
{\small\setlength{\tabcolsep}{3pt}
\begin{tabularx}{\textwidth}{XX}
\hline
Device & Architecture
\\\hline
D1 & FTO / PCBM / RbGeI$_{3}$ / NiO / Au\\
D2 & FTO / TiO$_{2}$ / RbGeI$_{3}$ / NiO / Au\\
D3 & FTO / PCBM / RbGeI$_{3}$ / CuI / Au\\
D4 & FTO / TiO$_{2}$ / RbGeI$_{3}$ / CuI / Au\\
D5 & FTO / PCBM / RbGeI$_{3}$ / CBTS / Au\\
D6 & FTO / TiO$_{2}$ / RbGeI$_{3}$ / CBTS / Au\\
D7 & FTO / PCBM / RbGeI$_{3}$ / Spiro-OMeTAD / Au\\
D8 & FTO / TiO$_{2}$ / RbGeI$_{3}$ / Spiro-OMeTAD / Au\\
\hline
\end{tabularx}
}
\caption{Eight device configurations screened in this study. All configurations were simulated at an absorber thickness of 400 nm.}
\end{table}

\section{Results and Discussion}

\subsection{Screening of device configurations}
Table II compiles the photovoltaics of all eight configurations: two ETLs (TiO$_{2}$, PCBM) crossed with four HTLs (NiO, CuI, CBTS, Spiro-OMeTAD) on the common RbGeI$_{3}$ absorber.
The D4 architecture (TiO$_{2}$/CuI) combines well-matched band alignment with an all-inorganic transport stack, avoiding the stability limitations of organic HTLs.

\begin{table}[htbp]
\centering
{\footnotesize\setlength{\tabcolsep}{3pt}
\begin{tabularx}{\textwidth}{XXcccc}
\hline
Dev. & Architecture & V$_{\mathrm{OC}}$ (V) & J$_{\mathrm{SC}}$ (mA/cm$^{2}$) & FF (\%) & PCE (\%)
\\\hline
D1 & FTO/\allowbreak{}PCBM/\allowbreak{}RbGeI$_{3}$/NiO/\allowbreak{}Au & 0.81 & 30.10 & 83.68 & 20.40\\
D2 & FTO/\allowbreak{}TiO$_{2}$/RbGeI$_{3}$/NiO/\allowbreak{}Au & 0.82 & 30.46 & 84.48 & 21.10\\
D3 & FTO/\allowbreak{}PCBM/\allowbreak{}RbGeI$_{3}$/CuI/\allowbreak{}Au & 0.81 & 30.10 & 83.96 & 20.47\\
D4 & FTO/\allowbreak{}TiO$_{2}$/RbGeI$_{3}$/CuI/\allowbreak{}Au & 0.81 & 30.46 & 79.98 & 19.78\\
D5 & FTO/\allowbreak{}PCBM/\allowbreak{}RbGeI$_{3}$/CBTS/\allowbreak{}Au & 0.82 & 30.20 & 82.36 & 20.40\\
D6 & FTO/\allowbreak{}TiO$_{2}$/RbGeI$_{3}$/CBTS/\allowbreak{}Au & 0.82 & 30.55 & 82.13 & 20.57\\
D7 & FTO/\allowbreak{}PCBM/\allowbreak{}RbGeI$_{3}$/Spiro/Au & 0.81 & 30.07 & 85.01 & 20.71\\
D8 & FTO/\allowbreak{}TiO$_{2}$/RbGeI$_{3}$/Spiro/Au & 0.82 & 30.46 & 83.95 & 20.97\\
\hline
\end{tabularx}
}
\caption{Initial photovoltaic performance of the eight screened device configurations, simulated at an absorber thickness of 400 nm.  Selected as the baseline for full parametric optimization.}
\end{table}
The best FF belongs to D7 (85.01\%) and the best PCE to D2 (21.10\%), not D4. Because D2 led at screening, it was run through the identical eleven-step pipeline on its own NiO deck (literature-standard NiO parameters, documented in S7; same triple-gate adoptions) as a like-for-like check: D2 carries 0.8~$\mu$m / $10^{12}$ defects / $\epsilon_{r}$ 23.1 / $\chi$ 3.9 / 10~nm ETL / $10^{18}$ densities to a raw 32.87\% chain end, and its gate-adopted device reaches 28.16\% (1.11~V, 30.32~mA/cm$^{2}$, 83.57\%) --- the screening gap survives optimization essentially intact (+1.32~pp at screening, +1.10~pp adopted), so D4 is not the efficiency optimum of this family. D4 is retained on stability and cost grounds stated plainly: CuI's documented stability and solution-processability advantages over Spiro-OMeTAD \cite{r25,r26} plus NiO's moisture sensitivity against D2, and an all-inorganic transport pair either way. Per-step D2 nodes are tabulated in Supplementary Section~S1. Screening ran once per configuration on the pre-rebase deck, so Table II is an architecture-ranking record, not headline numbers; the selected D4 stack was fully re-optimized on the rebased frozen deck. Band-alignment conclusions (offsets, extraction paths) transfer across the rebase while absolute PCEs do not.

\subsection{Optimization summary}
Eleven parameters were optimized one at a time on the screened
FTO/TiO$_{2}$/RbGeI$_{3}$/CuI/\allowbreak{}Au stack (absorber thickness, bulk
and interface defect densities, dielectric constant, electron affinity,
bandgap, transport-layer thicknesses, conduction- and valence-band densities
of states); each value was carried forward under the triple-gate rule, giving the consolidated set in
Table IV. Full per-parameter traces (Figs.\ S1--S10) are given in
Supplementary Section~S1, and the illumination and dark characteristics
(Figs.\ S11--S12) in Section~S2. The defect-density step itself is retained
below (Section~V.C) together with its energy/asymmetry extension, since
Ge-related degradation is central to this work.

\begin{table}[htbp]
\centering
{\footnotesize\setlength{\tabcolsep}{3pt}
\begin{tabularx}{\textwidth}{Xccc}
\hline
Step & Sweep range & Adopted & Gate basis
\\\hline
1 thickness & 0.3--1.0 $\mu$m & 0.7 $\mu$m & sweep max, plateau \\
2 bulk $N_{t}$ & 1e12--1e18 cm$^{-3}$ & 1e14 & passivated-film pin (max at 1e12 shown) \\
3 $\epsilon_{r}$ & 10--35 & 15 & sweep-flat pin \\
4 $\chi$ & 3.7--4.2 eV & 3.9 eV & interior max (3.7 non-convergent) \\
5/6 interface $N_{t}$ & 1e12--1e18 cm$^{-2}$ & 1e14 & passivated-film pin (max at 1e12 shown) \\
7 $E_{g}$ & 1.2--1.7 eV & 1.4 eV & SQ-margin gate (max at 1.3 rejected) \\
8 TiO$_{2}$ & 10--90 nm & 10 nm & sweep max + pinhole caveat \\
9 CuI & 0.05--1.9 $\mu$m & 100 nm & sweep-flat pin \\
10/11 $N_{C}$/$N_{V}$ & 1e18--1e20 cm$^{-3}$ & 1e18 & literature floor (sub-floor $\le$0.13 pp) \\
\hline
\end{tabularx}
}
\caption{Per-step sweeps, adopted values and gate basis. Full traces: Figs.~S1--S10.}
\end{table}
\subsection{Effect of absorber defect density}
Swept $10^{12}$--$10^{18}$ cm$^{-3}$, PCE peaks at 25.27\% at the $10^{12}$ cm$^{-3}$ floor, reads 23.37\% at the adopted passivated-film pin of $10^{14}$ cm$^{-3}$, then collapses (23.37 $\rightarrow$ 6.79\%) through falling V$_{\mathrm{OC}}$ (0.964 $\rightarrow$ 0.559 V), saturating past $10^{17}$ cm$^{-3}$.
Beyond the total trap density, the energetic position $E_{\mathrm{t}}$ of the dominant defect and the electron/hole capture asymmetry ($\sigma_{n}$, $\sigma_{p}$) were swept to model Ge-related degradation explicitly. At the H1 baseline ($N_{\mathrm{t}}$ = $10^{14}$ cm$^{-3}$), moving $E_{\mathrm{t}}$ from 0.2 eV to 1.2 eV above $E_{\mathrm{v}}$ traces a shallow symmetric minimum centred at mid-gap: PCE falls from 27.27\% at the band edges to 27.06\% for $E_{\mathrm{t}}$ = 0.4-1.0 eV (Fig. 10), the textbook SRH fingerprint that mid-gap states are the most lethal recombination centres. The level position is thus immaterial at baseline densities (27.06-27.27\%). A two-point check at $N_{\mathrm{t}}$ = $10^{15}$ cm$^{-3}$ gives 25.07\% (0.2 eV) and 24.40\% (0.7 eV), consistent with the bulk-ladder penalty rather than a level effect. Capture asymmetry at mid-gap (Fig. 11) is significant: symmetric large traps ($\sigma_{n}$=$\sigma_{p}$=$10^{-17}$ m$^{2}$) collapse to 20.76\%, the worst case; at fixed large $\sigma_{n}$=$10^{-17}$ m$^{2}$, shrinking $\sigma_{p}$ recovers 21.91\% ($10^{-19}$ m$^{2}$) and then 27.59\% ($10^{-21}$ m$^{2}$); small $\sigma_{n}$ ($\le$$10^{-19}$ m$^{2}$) is benign at any $\sigma_{p}$ (26.99--27.80\%). Electron capture sets the damage. The oxidation threat (Ge$^{2+}$ $\rightarrow$ Ge$^{4+}$, Ge vacancies) therefore cannot be reduced to a single $N_{\mathrm{t}}$ number: its energy and capture character jointly set the damage, and dense symmetric traps are the worst case.

\begin{figure}[htbp]
\centering
\includegraphics[width=0.85\linewidth]{figs/fig4.png}
\caption{Variation of PCE, V$_{\mathrm{OC}}$, J$_{\mathrm{SC}}$, and FF with RbGeI$_{3}$ absorber defect density. Performance degrades sharply above 10$^{14}$ cm$^{-3}$.}
\end{figure}

\subsection{Energy band diagram analysis}
The band diagram under illumination (Figure 3) shows clean extraction at both heterojunctions. At CuI/RbGeI$_{3}$ the 1.8 eV conduction offset blocks electrons while the 0.10 eV valence step passes holes; at RbGeI$_{3}$/TiO$_{2}$ a 0.10 eV conduction step assists extraction while a 1.90 eV valence offset blocks holes. The peak bulk quasi-Fermi splitting reaches $\approx$1.13 V while the terminal V$_{\mathrm{OC}}$ is 1.06 V: the 70 mV gap is interface recombination plus quasi-Fermi bending toward the contacts (visible in Fig.~3), i.e.\ a quantified contact loss, not an inconsistency.

\begin{figure}[htbp]
\centering
\includegraphics[width=0.85\linewidth]{figs/fig15.png}
\caption{Energy band diagram of the optimized FTO/\allowbreak{}TiO$_{2}$/RbGeI$_{3}$/CuI/\allowbreak{}Au solar cell under illumination, showing the conduction-band edge (EC), valence-band edge (EV), electron quasi-Fermi level ($F_{\mathrm{n}}$), and hole quasi-Fermi level ($F_{\mathrm{p}}$). The 1.8 eV conduction-band offset at the CuI/\allowbreak{}RbGeI$_{3}$ interface suppresses electron back-transfer, while the 0.10 eV conduction-band offset at the RbGeI$_{3}$/TiO$_{2}$ interface aids electron extraction. The 1.90 eV valence-band offset at the RbGeI$_{3}$/TiO$_{2}$ interface provides an energetic barrier against hole injection into TiO$_{2}$.}
\end{figure}

\subsection{Current density--voltage characteristics}
The consolidated J--V under AM 1.5G (Figure 4) is sharply rectangular: V$_{\mathrm{OC}}$ = 1.06 V, J$_{\mathrm{SC}}$ = 30.32 mA/cm$^{2}$, FF = 84.11\%, PCE = 27.06\%, with no s-kink or rollover.

\begin{figure}[htbp]
\centering
\includegraphics[width=0.85\linewidth]{figs/fig16.png}
\caption{Current density--voltage (J--V) curve of the optimized FTO/\allowbreak{}TiO$_{2}$/RbGeI$_{3}$/CuI/\allowbreak{}Au device under AM 1.5G illumination. The red dashed curve shows the output power density, with the maximum-power point marked.}
\end{figure}

\subsection{Summary of optimized device parameters and final performance}
After completing the eleven sequential steps, the adopted values were incorporated into the proposed FTO/\allowbreak{}TiO$_{2}$/RbGeI$_{3}$/CuI/\allowbreak{}Au cell (Table IV): adopted means triple-gate values (Section~IV.D), not raw sweep maxima --- defect densities pinned at the passivated-film $10^{14}$ level, N$_{\mathrm{C}}$/N$_{\mathrm{V}}$ at the $10^{18}$ literature floor, dielectric constant at 15 (sweep-flat), CuI at 100 nm (sweep-flat). The consolidated device was simulated in a single SCAPS-1D run with detailed-balance-consistent radiative recombination throughout: V$_{\mathrm{OC}}$ = 1.06 V, J$_{\mathrm{SC}}$ = 30.32 mA/cm$^{2}$, FF = 84.11\%, PCE = 27.06\%. The V$_{\mathrm{OC}}$ climb from the 0.81 V screening point is accounted for: bulk $N_{t}$ $10^{15}$$\rightarrow$$10^{14}$ adds $\approx$0.09 V (0.877$\rightarrow$0.964 V in the sweep), and the $N_{C}$/$N_{V}$ literature floor adds $\approx$0.09 + 0.06 V across their sweeps, with the remainder from thickness/affinity alignment --- no single step carries it. Every metric sits below its R\"uhle Shockley--Queisser limit at 1.4 eV (32.88 mA/cm$^{2}$, 1.122 V, 89.3\%, 32.91\%) \cite{r36}. A six-parameter genetic-algorithm search (population 10, six generations, 62 SCAPS evaluations; crossover plus Gaussian mutation with elitism; seeded, archived) returns a best-found joint point at $E_{\mathrm{g}}$ = 1.36 eV, $N_{\mathrm{t}}$ $\approx$ 1.6 $\times$ 10$^{13}$ cm$^{-3}$, interface $N_{\mathrm{t}}$ $\approx$ 6 $\times$ 10$^{11}$ cm$^{-2}$, N$_{\mathrm{C}}$/N$_{\mathrm{V}}$ $\approx$ 1.6 $\times$ 10$^{18}$ cm$^{-3}$, absorber 0.72 $\mu$m, giving 1.04 V, 31.55 mA/cm$^{2}$, 85.87\%, 28.30\% (Fig. 7, Table V) --- also fully SQ-admissible, but reported as best-found optimizer reach rather than an admissible device: its $N_{t}$$\approx$1.6$\times$10$^{13}$ cm$^{-3}$ and interface $N_{t}$$\approx$6$\times$10$^{11}$ cm$^{-2}$ sit below the triple-gate floors, not a proven global maximum. Section VI.C stacks all realism losses onto Table IV for a single realistic window (Fig. 9).

\begin{table}[htbp]
\centering
{\footnotesize\setlength{\tabcolsep}{3pt}
\begin{tabularx}{\textwidth}{Xcccc}
\hline
Parameter & TiO$_{2}$ & FTO & RbGeI$_{3}$ & CuI
\\\hline
Thickness (nm) & 10 & 500 & 700 & 100\\
Bandgap $E_{\mathrm{g}}$ (eV) & 3.2 & 3.2 & 1.4 & 3.1\\
Electron affinity $\chi$ (eV) & 4.0 & 4.4 & 3.9 & 2.1\\
Dielectric permittivity $\epsilon$r & 9 & 9 & 15 & 6.5\\
CB effective DOS N$_{\mathrm{C}}$ (cm$^{-3}$) & 2$\times$10$^{18}$ & 2.2$\times$10$^{18}$ & 1$\times$10$^{18}$ & 2.8$\times$10$^{19}$\\
VB effective DOS N$_{\mathrm{V}}$ (cm$^{-3}$) & 1.8$\times$10$^{19}$ & 1.8$\times$10$^{19}$ & 1$\times$10$^{18}$ & 1$\times$10$^{19}$\\
Electron mobility $\mu$e (cm$^{2}$ V$^{-1}$ s$^{-1}$) & 20 & 20 & 28.6 & 100\\
Hole mobility $\mu$p (cm$^{2}$ V$^{-1}$ s$^{-1}$) & 10 & 10 & 27.3 & 43.9\\
Donor density $N_{\mathrm{D}}$ (cm$^{-3}$) & 9$\times$10$^{16}$ & 1$\times$10$^{19}$ & 1$\times$10$^{9}$ & 0\\
Acceptor density $N_{\mathrm{A}}$ (cm$^{-3}$) & 0 & 0 & 1$\times$10$^{9}$ & 1$\times$10$^{18}$\\
Bulk defect density $N_{\mathrm{t}}$ (cm$^{-3}$) & 1$\times$10$^{15}$ & 1$\times$10$^{14}$ & 1$\times$10$^{14}$ & 1$\times$10$^{14}$\\
\hline
\end{tabularx}
}
\caption{Optimized material parameters of the proposed FTO/\allowbreak{}TiO$_{2}$/RbGeI$_{3}$/CuI/\allowbreak{}Au perovskite solar cell. Transport-layer bulk defect densities were carried from the initial deck (not among the eleven swept steps).}
\end{table}
\textit{Note: The optimized interface defect density was 1 $\times$ 10$^{14}$ cm$^{-2}$ for both the TiO}\textit{2}\textit{/RbGeI}\textit{3}\textit{ and RbGeI}\textit{3}\textit{/CuI interfaces.}

\begin{table}[htbp]
\centering
{\footnotesize\setlength{\tabcolsep}{3pt}
\begin{tabularx}{\textwidth}{Xcccc}
\hline
Device & V$_{\mathrm{OC}}$ (V) & J$_{\mathrm{SC}}$ (mA/cm$^{2}$) & FF (\%) & PCE (\%)
\\\hline
Adopted reference (Table IV) & 1.06 & 30.32 & 84.11 & 27.06\\
GA best-found ($E_{\mathrm{g}}$ 1.36 eV) & 1.04 & 31.55 & 85.87 & 28.30\\
Reference, N$_{\mathrm{C}}$=N$_{\mathrm{V}}$=1e19 & 0.97 & 30.32 & 84.09 & 24.82\\
Reference + realistic optics (0.9 sun) & 1.06 & 27.28 & 83.97 & 24.28\\
\hline
\end{tabularx}
}
\caption{Adopted reference device versus GA best-found point and sensitivity rows (defect rows re-simulated, Fig. 8; optics row reports PCE$_{1\text{-sun}}$). PCEs are unrounded SCAPS outputs (displayed inputs rounded).}
\end{table}
As an arithmetic consistency check, the final-device PCE computed from Eq.~(\ref{eq:pce}) using the consolidated J--V metrics is PCE = (1.0613 $\times$ 30.3161 $\times$ 0.8411) / 100 = 27.06\%, matching the SCAPS-1D output to within rounding.

\subsection{Uncertainty quantification}
Monte Carlo UQ used Latin hypercube sampling (N = 200, seed 11) over eight variables on the H1 deck: gap 1.35--1.45 eV ($E_{g}$ moves co-move $\chi$), bulk $5\times10^{13}$--$5\times10^{15}$ cm$^{-3}$, interfaces $5\times10^{13}$--$5\times10^{14}$ cm$^{-2}$ each side, $N_{\mathrm{C}}$/$N_{\mathrm{V}}$ $5\times10^{16}$--$5\times10^{17}$ cm$^{-3}$, $\mu_{n}$/$\mu_{p}$ 15--50 cm$^{2}$/Vs.

\begin{figure}[htbp]
\centering
\includegraphics[width=0.85\linewidth]{figs/fig17.png}
\caption{Uncertainty quantification results (N = 200). Histograms of PCE, V$_{\mathrm{OC}}$, J$_{\mathrm{SC}}$, and FF.}
\end{figure}
Mean PCE 24.86 $\pm$ 1.77 \% (median 24.55 \%; 5th/95th 22.23/27.79 \%; 200/200 converged).

\subsection{Temperature dependence}
Swept 280--400 K (20 K steps) under AM 1.5G.

\begin{figure}[htbp]
\centering
\includegraphics[width=0.85\linewidth]{figs/fig20.png}
\caption{Temperature dependence of photovoltaic parameters from 280 K to 400 K.}
\end{figure}
dV$_{\mathrm{OC}}$/dT = $-$1.38 mV/K, dPCE/dT = $-$0.058 \%/K (PCE falls non-monotonically in FF, so the linear coefficient is descriptive, not predictive); $\sim$75 \% of RT PCE retained at 400 K. V$_{\mathrm{OC}}$(T) extrapolates to $E_{\mathrm{a}}$ = 1.48 eV --- 0.08 eV \textit{above} $E_{\mathrm{g}}$, so this is not a clean bulk-recombination signature but an extrapolation artifact consistent with T-dependent N$_{\mathrm{C}}$/N$_{\mathrm{V}}$; it is reported as measured, with the artifact disclosed, rather than as proof of bulk-dominated recombination.

\begin{figure}[htbp]
\centering
\includegraphics[width=0.85\linewidth]{figs/fig21.png}
\caption{Genetic-algorithm joint search over bandgap, bulk/interface defect densities, N$_{\mathrm{C}}$, N$_{\mathrm{V}}$, and absorber thickness. Per-generation best (grey) fluctuates under mutation; best-so-far (squares) reaches 28.30\% at generation 5 from 27.61\% at generation 0.}
\end{figure}

\begin{figure}[htbp]
\centering
\includegraphics[width=0.85\linewidth]{figs/fig22.png}
\caption{Defect tolerance of the adopted device. Raising bulk $N_{\mathrm{t}}$ to $10^{15}$ ($10^{16}$) cm$^{-3}$ drops PCE to 24.40\% (20.76\%); raising rear-interface $N_{\mathrm{t}}$ to $10^{15}$ ($10^{16}$) cm$^{-2}$ drops PCE to 20.29\% (19.65\%). J$_{\mathrm{SC}}$ is nearly unchanged; V$_{\mathrm{OC}}$ and FF collapse.}
\end{figure}

\begin{figure}[htbp]
\centering
\includegraphics[width=0.85\linewidth]{figs/fig23.png}
\caption{From the adopted 27.06\% to the stacked realistic PCE: optics, bulk and interface losses give 20.78\% (21.76\% with BSF). All nodes re-simulated PCE$_{1\text{-sun}}$ on the adopted deck.}
\end{figure}
\subsection{Effect of background acceptor doping}
Each Ge$^{2+}$\,$\rightarrow$\,Ge$^{4+}$ oxidation event donates two holes via V$_{\mathrm{Ge}}$, so background doping is oxidation by another name: sweeping $N_{\mathrm{A}}$ $10^{14}$--$10^{17}$ cm$^{-3}$ on the stacked point (Fig.~12), PCE peaks at 23.41\% near $3\times10^{15}$ cm$^{-3}$ and softens past $10^{16}$ cm$^{-3}$ (23.12\%) before falling at higher doping (22.17\% at $10^{17}$) --- consistent with the oxidation budget's $p\approx2$[V$_{\mathrm{Ge}}$] expectation in Section~VI.H, with the penalty soft ($\le$0.3 pp at $10^{16}$), not a collapse. Synthetic target: $N_{\mathrm{A}}\le10^{16}$ cm$^{-3}$. Tying doping to defects per the budget ($N_{t}$=$10^{15}$ with $N_{A}$=$2\times10^{15}$ cm$^{-3}$) gives 25.52\% against 24.40\% untied, so the near-intrinsic absorber used throughout is conservative at this rung, not flattering.
\subsection{Effect of series resistance: manufacturing tolerance}
Corrections were recomputed on the stacked realistic point (20.78\%), not the ideal: FF falls 77.75 $\rightarrow$ 59.75\% and PCE 20.78 $\rightarrow$ 15.96\% as $R_{\mathrm{s}}$ goes 0 $\rightarrow$ 8 $\Omega\cdot$cm$^{2}$ (Fig.~13, external V $-$ J$\times$$R_{\mathrm{s}}$). Tolerance: $R_{\mathrm{s}}\le4$ preserves 20\% pre-optics (23.16 $\rightarrow$ 20.10\%); on the full stack the same $R_{\mathrm{s}}$ gives 18.31\%. Shunt analysis on the same point: $R_{\mathrm{sh}}$ $10^{2}$/$5\times10^{2}$/$10^{3}$/$\infty$ gives 14.26/19.41/20.09/20.71\%, so $R_{\mathrm{sh}}\ge10^{3}$ $\Omega\cdot$cm$^{2}$ holds 20\%.
\subsection{Absorber mobility viability}
Electron/hole mobilities were swept jointly from 1 to 100 cm$^{2}$/Vs on the stacked point (Fig.~14): PCE is log-linear (18.35\% at 1 to 24.32\% at 100). Crystallization target: $\mu\ge15$ cm$^{2}$/Vs.


\section{Extended results: loss budget, sensitivity and experimental grounding}
\subsection{Benchmarking data (discussion in Section~VII)}
Literature benchmarking (Table~VI) and the CsGeI$_{3}$ gap control (Table~VII, Fig.~16) are discussed in Section~VII; the tables and figures are placed here.

\begin{table}[htbp]
\centering
{\scriptsize\setlength{\tabcolsep}{2pt}
\begin{tabularx}{\textwidth}{XccccX}
\hline
Device Structure & V$_{\mathrm{OC}}$ (V) & J$_{\mathrm{SC}}$ (mA/cm$^{2}$) & FF (\%) & PCE (\%) & Year [Ref.]
\\\hline
FTO/\allowbreak{}TiO$_{2}$/RbGeI$_{3}$/NiO/\allowbreak{}Ag & 0.5311 & 28.89 & 63.68 & 9.77 & 2022 \cite{r14}\\
FTO/\allowbreak{}TiO$_{2}$/CsGeI$_{3}$/CuI/\allowbreak{}Au & 0.667 & 22.06 & 73.49 & 10.80 & 2022 \cite{r24}\\
ITO/\allowbreak{}C$_{60}$/RbGeI$_{3}$/CBTS/\allowbreak{}Ag & 0.99 & 33.20 & 82.80 & 27.21 & 2024 \cite{r22}\\
FTO/\allowbreak{}TiO$_{2}$/RbGeI$_{3}$/PEDOT:PSS/\allowbreak{}Au & 0.813 & 30.84 & 74.19 & 18.60 & 2025 \cite{r39}\\
FTO/\allowbreak{}TiO$_{2}$/MASnI$_{3}$/CBTS/\allowbreak{}Au & 1.030 & 30.24 & 83.41 & 25.98 & 2026 \cite{r41}\\
FTO/\allowbreak{}TiO$_{2}$/MASnI$_{3}$/CZTS/Au & 0.96 & 31.66 & 67.00 & 20.36 & 2025 \cite{r42}\\
FTO/\allowbreak{}TiO$_{2}$/CsSnI$_{3}$/P3HT/Au & 0.972 & 34.70 & 78.21 & 26.38 & 2024 \cite{r40}\\
FTO/\allowbreak{}TiO$_{2}$/RbGeI$_{3}$/Spiro-OMeTAD/\allowbreak{}Au & 1.1813 & 26.31 & 85.16 & 26.47 & 2026 \cite{r23}\\
FTO/\allowbreak{}TiO$_{2}$/RbGeI$_{3}$/CuI/\allowbreak{}Au (this work) & 1.06 & 30.32 & 84.11 & 27.06 & This work\\
FTO/\allowbreak{}TiO$_{2}$/RbGeI$_{3}$/Cu$_{2}$O-BSF/CuI/\allowbreak{}Au (this work) & 1.09 & 30.32 & 84.62 & 28.01 & This work\\
\hline
\end{tabularx}
}
\caption{Comparison of the proposed device with recently reported lead-free perovskite solar cells. All comparator values verified against primary sources. All entries are SCAPS-1D simulation studies.}
\end{table}

\begin{table}[htbp]
\centering
{\footnotesize\setlength{\tabcolsep}{3pt}
\begin{tabularx}{\textwidth}{XXccc}
\hline
Device & Condition & V$_{\mathrm{OC}}$ (V) & J$_{\mathrm{SC}}$ (mA/cm$^{2}$) & PCE (\%)
\\\hline
RbGeI$_{3}$ & Ideal ($N_{\mathrm{t}}$=1e14) & 1.06 & 30.32 & 27.06\\
CsGeI$_{3}$ & Ideal ($N_{\mathrm{t}}$=1e14) & 1.10 & 23.18 & 19.65\\
RbGeI$_{3}$ & Realistic ($N_{\mathrm{t}}$=1e15) & 1.02 & 30.31 & 24.40\\
CsGeI$_{3}$ & Realistic ($N_{\mathrm{t}}$=1e15) & 1.10 & 23.17 & 19.31\\
RbGeI$_{3}$ & High oxidation ($N_{\mathrm{t}}$=1e16) & 0.96 & 30.22 & 20.76\\
CsGeI$_{3}$ & High oxidation ($N_{\mathrm{t}}$=1e16) & 1.09 & 23.10 & 17.07\\
\hline
\end{tabularx}
}
\caption{RbGeI$_{3}$ versus CsGeI$_{3}$ in the identical FTO/\allowbreak{}TiO$_{2}$/absorber/CuI/\allowbreak{}Au stack under ideal, realistic, and high-oxidation defect budgets. Only the absorber bandgap differs (1.4 vs 1.6 eV) at fixed affinity --- a gap-effect control, not A-site chemistry proof.}
\end{table}

\begin{figure}[htbp]
\centering
\includegraphics[width=0.85\linewidth]{figs/fig24.png}
\caption{PCE versus defect energy $E_{\mathrm{t}}$ on the H1 baseline ($N_{\mathrm{t}}$ = $10^{14}$ cm$^{-3}$). Symmetric minimum at mid-gap (27.06\%) recovering to 27.27\% at both band edges: the SRH fingerprint.}
\end{figure}

\begin{figure}[htbp]
\centering
\includegraphics[width=0.85\linewidth]{figs/fig25.png}
\caption{PCE versus capture cross-section asymmetry at mid-gap $E_{\mathrm{t}}$ on the H1 baseline. Symmetric large traps (20.76\%) are worst; at $\sigma_{n}$=$10^{-17}$ m$^{2}$, shrinking $\sigma_{p}$ recovers 21.91\% then 27.59\%; small $\sigma_{n}$ is benign throughout.}
\end{figure}

\begin{figure}[htbp]
\centering
\includegraphics[width=0.85\linewidth]{figs/fig26.png}
\caption{PCE versus background acceptor density on the stacked point. Optimum near $3\times10^{15}$ cm$^{-3}$ (23.41\%); penalty soft at $10^{16}$ cm$^{-3}$ ($\le$0.3 pp), collapse beyond.}
\end{figure}

\begin{figure}[htbp]
\centering
\includegraphics[width=0.85\linewidth]{figs/fig27.png}
\caption{PCE and FF versus series resistance (external V - J x $R_{\mathrm{s}}$ correction) on the stacked 20.78\% point. PCE falls to 15.96\% at $R_{\mathrm{s}}$ = 8 $\Omega\cdot$cm$^{2}$; $R_{\mathrm{s}}\le4$ preserves 20\% only pre-optics.}
\end{figure}

\begin{figure}[htbp]
\centering
\includegraphics[width=0.85\linewidth]{figs/fig28.png}
\caption{PCE versus absorber mobility ($\mu_{\mathrm{n}}$ = $\mu_{\mathrm{p}}$) on the stacked point. Log-linear rise 18.35\% (1) to 24.32\% (100 cm$^{2}$/Vs); target $\mu$ >= 15 cm$^{2}$/Vs.}
\end{figure}

\begin{figure}[htbp]
\centering
\includegraphics[width=0.85\linewidth]{figs/fig29.png}
\caption{V$_{\mathrm{OC}}$ versus temperature with 0 K extrapolation. $E_{\mathrm{a}}$ = 1.48 eV sits 0.08 eV \textit{above} $E_{\mathrm{g}}$: extrapolation artifact (T-dependent $N_{C}$/$N_{V}$), not a bulk-recombination signature; dV$_{\mathrm{OC}}$/dT = -1.38 mV/K.}
\end{figure}

\begin{figure}[htbp]
\centering
\includegraphics[width=0.85\linewidth]{figs/fig30.png}
\caption{RbGeI$_{3}$ versus CsGeI$_{3}$ under ideal, realistic, and high-oxidation defect budgets in the identical stack. RbGeI$_{3}$ leads on PCE at every level (gap harvest); Cs V$_{\mathrm{OC}}$ pinning signals interface limitation.}
\end{figure}

\begin{figure}[htbp]
\centering
\includegraphics[width=0.85\linewidth]{figs/fig31.png}
\caption{Resistance derating ladder on the stacked 20.78\% point ($R_{\mathrm{s}}$/$R_{\mathrm{sh}}$ post-corrections; red dashed line marks 20\%).}
\end{figure}

\subsection{BSF enhancement of the TiO$_{2}$/RbGeI$_{3}$/CuI stack}
A Cu$_{2}$O-like p$^{+}$ back-surface field (BSF) inserted between RbGeI$_{3}$ and CuI ($E_{\mathrm{g}}$ 2.1 eV, $\chi$ 3.2 eV, $\epsilon_{r}$ 7.5, literature-standard Cu$_{2}$O parameters for SCAPS perovskite studies \cite{rMA}) upgrades the same economical stack without changing the absorber or transport pair. The 50 nm, $N_{A}$ $10^{18}$ cm$^{-3}$ pilot converges to 27.93\% (+0.87 pp over the 27.06\% headline). The full $4\times3$ grid (Fig.~18; nodes in Supplementary Table~S5) shows doping, not thickness, controls the gain: $10^{17}$ gives 27.55--27.60\%, $10^{18}$ gives 27.93--27.94\%, $10^{19}$ gives 28.00--28.01\%. The optimum reaches 28.01\% (1.0915 V, 30.32 mA/cm$^{2}$, 84.62\%); $J_{\mathrm{SC}}$ is unchanged, so the gain is $V_{\mathrm{OC}}$ (+30 mV) plus FF (+0.5 pp). The mechanism is rear-interface gap widening, not electron reflection (CuI already blocks electrons with 1.8 eV): the BSF removes the 0.10 eV valence cliff -- computed from the stated layer offsets, cf.\ the no-BSF lineup in Fig.~3 -- that otherwise shrinks the interface recombination gap, and it halves the rear trap penalty --- raising rear $N_{t}$ to $10^{16}$ cm$^{-2}$ costs 7.41 pp without BSF (27.06 $\rightarrow$ 19.65\%) but 3.64 pp with it (28.00 $\rightarrow$ 24.36\%, Fig.~18, right panel). Block attribution is explicit: both rear blocks (CuI/BSF and BSF/RbGeI$_{3}$) carry the raised $N_{t}$, and raising only the far CuI/BSF block costs nothing, pinning the damage on the absorber side. A p$^{+}$ control interlayer with CuI bands (no new offset, same $N_{A}$) gains only +0.76 pp ideally (27.82\%) and still loses 6.49 pp under traps (21.33\%), so the field contributes a minority share and the offset removal carries the immunity. Gains are non-additive with the $N_{t}$=$10^{15}$ ladder for the same reason: once the rear gap is widened, further defect cuts buy less. All metrics sit below the R\"uhle limits \cite{r35,r36,r37}. Cu$_{2}$O is routinely sputtered on perovskite stacks, but a direct Cu$_{2}$O-on-Ge-perovskite interface is undemonstrated and flagged as the fabrication risk.

\begin{figure}[htbp]
\centering
\includegraphics[width=0.85\linewidth]{figs/fig32.png}
\caption{BSF grid PCE (\%): $N_{A}$ $10^{17}\!\rightarrow\!10^{19}$ cm$^{-3}$ rows; thickness 30--100 nm columns. Doping dominates; thickness is flat. Full node values: Supplementary Table~S5. Right: rear-interface trap attribution --- raising rear $N_{t}$ to $10^{16}$ cm$^{-2}$ costs 7.41 pp without BSF (27.06 $\rightarrow$ 19.65\%) but 3.64 pp with it (28.00 $\rightarrow$ 24.36\%); far-side traps alone cost nothing; the p+ no-offset control (grey) traps at 21.33\%.}
\end{figure}

\subsection{Loss waterfall: what survives realism}
Native-defect simulation plus photon-scaling optics and post-processed parasitics bound the experimentally plausible window (Fig.~19; nodes in Supplementary Table~S6; PCE$_{1\text{-sun}}$ norm.). Radiative recombination at $K_{\mathrm{rad}}$ $10^{-10}$ cm$^{3}$/s costs 0.07 pp (27.13 $\rightarrow$ 27.06\%), consistent with the detailed-balance value to order of magnitude at the rebased densities; higher $K_{\mathrm{rad}}$ bounds the sensitivity in S6. Auger coefficients $10^{-29}$--$10^{-27}$ cm$^{6}$/s change nothing at 1-sun ($\pm$0.00 pp) --- a reportable null justifying its omission from the headline budget. Ten percent front reflection costs 2.78 pp ($\rightarrow$24.28\%); realistic bulk defects ($10^{15}$ cm$^{-3}$) a further 2.35 pp ($\rightarrow$21.93\%); realistic interfaces ($10^{15}$ cm$^{-2}$) a further 1.15 pp ($\rightarrow$\textbf{20.78\%}). The BSF optimum on the \textit{identical} full stack gives \textbf{21.76\%}, recovering +0.98 pp. The stack is optics- and defect-limited, not transport-limited, and the BSF pays off precisely where reality lives.

\begin{figure}[htbp]
\centering
\includegraphics[width=0.85\linewidth]{figs/fig23.png}
\caption{Stacked loss waterfall (PCE$_{1\text{-sun}}$ norm.): optics $-$2.78 pp, bulk $-$2.35 pp, interfaces $-$1.15 pp; BSF on the full stack recovers +0.98 pp. Nodes: Supplementary Table~S6.}
\end{figure}

\subsection{Carrier-lifetime baseline}
SRH statistics ($\tau^{-1}$=$\sigma v_{\mathrm{th}}N_{t}$) give \textbf{1000 ns} at $10^{14}$ cm$^{-3}$ with diffusion length \textbf{8.6 $\mu$m} $\gg$ 700 nm absorber, explaining the near-unity collection. Fig.~20 shows the Beer--Lambert generation profile normalized to the H1 $J_{\mathrm{SC}}$. No built-in potential, depletion width or Mott--Schottky values are reported: the near-intrinsic absorber is fully depleted so the depletion approximation does not apply, and C--V extraction is left for fabricated devices (method note in S7).

\begin{figure}[htbp]
\centering
\includegraphics[width=0.85\linewidth]{figs/fig34.png}
\caption{Normalized Beer--Lambert generation profile (square-root law, edge 886 nm). Recombination-profile export is unavailable for this deck (S7); the BSF mechanism rests on the controlled attribution runs (Fig.~18) instead.}
\end{figure}

\subsection{Data-driven global sensitivity analysis}
An XGBoost surrogate on the 200-point LHS-UQ block (8 features: $E_{g}$, $N_{t}$, $N_{C}$, $N_{V}$, $\mu_{n}$, $\mu_{p}$, both interface $N_{t}$; 80/20 held-out split, seed 7; 200/200 converged; $E_{g}$ moves co-move $\chi$ so gap physics, not misalignment, is measured) gives held-out R$^{2}$ 0.907, RMSE 0.57 (Fig.~21; random forest agrees at 0.907). SHAP ranks bulk $N_{t}$ $>$ $E_{g}$ $>$ $\mu_{n}$ $>$ CuI-side $N_{t}$ $>$ TiO$_{2}$-side $N_{t}$ $>$ $N_{V}$ $>$ $\mu_{p}$ $>$ $N_{C}$ (Fig.~22): near the adopted point, bulk defect density dominates, electron mobility outranks hole mobility, and $N_{C}$/$N_{V}$ are inert \textit{within the sampled neighborhood} ($5\times10^{16}$--$5\times10^{17}$ cm$^{-3}$) --- consistent with Table V, where $10^{19}$ costs 2.2 pp outside it. Three hedges: rankings inherit sampled ranges ($E_{g}$ spans 1.35--1.45 while $N_{t}$ spans two decades); $N_{C}$/$N_{V}$ are sampled an order of magnitude below the enforced $10^{18}$ floor, so their ``inert'' ranking is a neighborhood statement, not a license to lower them; the interface range tops out at $5\times10^{14}$ cm$^{-2}$, excluding the realistic point. The deterministic-solver R$^{2}$ is reported on held-out data only; no factorial cross-check figure is shown (the prior 193-row block ran on the pre-rebase deck).

\begin{figure}[htbp]
\centering
\includegraphics[width=0.85\linewidth]{figs/fig35.png}
\caption{XGBoost parity on the UQ block, held-out 20\% (R$^{2}$ 0.907).}
\end{figure}

\begin{figure}[htbp]
\centering
\includegraphics[width=0.85\linewidth]{figs/fig36.png}
\caption{SHAP importance from the UQ surrogate (R$^{2}$ 0.907): bulk defects dominate, then gap and electron mobility.}
\end{figure}

\subsection{Pipeline calibration against an experimental inorganic perovskite}
No RbGeI$_{3}$ device has been fabricated to date, so direct validation is impossible; we instead run a calibration sanity check of the solver pipeline on the experimentally realized analog $\beta$-CsPbI$_{3}$ of Wang et al.\ \cite{rWang} (FTO/c-TiO$_{2}$/CsPbI$_{3}$/spiro-OMeTAD/Ag, 375 nm absorber, J--V measured at 318 K). Rebuilding that stack with literature-pinned bands and fitting only bulk $N_{t}$ $\times$ interface $N_{t}$ reproduces both reported devices (Table VIII, Fig.~23): the CHI-passivated champion at $N_{t}$ $3\times10^{13}$ cm$^{-3}$ / $N_{t,\mathrm{if}}$ $3\times10^{13}$ cm$^{-2}$ (18.49\% vs 18.4\%, all metrics within 2\%), and the control at $N_{t}$ $3\times10^{13}$ cm$^{-3}$ / $N_{t,\mathrm{if}}$ $10^{14}$ cm$^{-2}$ plus an effective $R_{s}$ $6.9$ $\Omega\cdot$cm$^{2}$ first-order correction (15.18\% vs 15.1\%). The control's low FF (0.72) cannot be reached by defects (floor $\sim$0.84) and carries the documented cracks, pinholes and hysteresis, justifying the resistive term; the control fit visibly misses the knee, so it is presented as achieved with that caveat. This calibrates that SCAPS-1D with literature bands can reproduce two measured inorganic devices by tuning defect densities (plus one resistive correction) --- it tests the solver pipeline, not CuI, Ge chemistry, the oxidation map or the BSF, none of which appear in the analog; full provenance is in Supplementary Section~S5.

\begin{table}[htbp]
\centering
{\footnotesize\setlength{\tabcolsep}{3pt}
\begin{tabularx}{\textwidth}{Xcccc}
\hline
Device & V$_{\mathrm{OC}}$ (V) & J$_{\mathrm{SC}}$ (mA/cm$^{2}$) & FF (\%) & PCE (\%)
\\\hline
Wang control (measured) & 1.05 & 20.03 & 72.0 & 15.10\\
SCAPS rebuild + $R_{s}$ 6.9 & 1.060 & 19.92 & 71.88 & 15.18\\
Wang champion (measured) & 1.11 & $\sim$20.0 & 82.0 & 18.40\\
SCAPS rebuild & 1.127 & 19.92 & 82.37 & 18.49\\
\hline
\end{tabularx}
}
\caption{Calibration check: SCAPS rebuild vs measured $\beta$-CsPbI$_{3}$ devices \cite{rWang} at 318 K. Fitted knobs are defect densities (plus $R_{s}$ on the control); the control fit misses the knee.}
\end{table}

\begin{figure}[htbp]
\centering
\includegraphics[width=\textwidth]{figs/fig37.png}
\caption{Calibration overlay: SCAPS rebuilds vs digitized J--V of Wang et al.\ \cite{rWang} control and CHI-passivated champion (318 K).}
\end{figure}

\subsection{Inorganic context from a frozen database snapshot}
We place the loss-budget window against the Perovskite Database Project via NOMAD \cite{rPDB} as era-bounded context, not a benchmark: 41,316 one-sun entries (43,106 raw fetch; validity filters in S5), of which 2,141 analysis-clean strictly inorganic entries (raw fetch 2,201) include 421 CsPbI$_{3}$ and just 5 Ge-bearing devices (zero RbGeI$_{3}$). The snapshot runs to early 2020 plus author uploads, so recent inorganic cells (20.6--21.5\%+) postdate it and no superiority over the present state of the art is claimed. Within the snapshot, our stacked 20.78\% point (21.76\% with BSF) sits above its inorganic max of 18.22\% (Fig.~24) --- a statement about the gap between simulation upper bounds and that era's inorganic state of the art, not evidence of realism; the refit champion marker (18.49\%) is our own calibration point plotted for scale, not an independent device, and its FF (82.37\%) sits just above the snapshot's realized FF max (0.82).

\begin{figure}[htbp]
\centering
\includegraphics[width=\textwidth]{figs/fig38.png}
\caption{(a) PCE histograms (NOMAD Perovskite DB \cite{rPDB}, frozen pre-2021 snapshot): all 1-sun, inorganic subset, CsPbI$_{3}$; gold band = this work's realistic window; dashed line = calibration refit. (b--d) V$_{\mathrm{OC}}$, FF, J$_{\mathrm{SC}}$ vs PCE; gold star = refit champion marker (18.49\%, plotted for scale), white star = control fit (15.18\%).}
\end{figure}

\subsection{Oxidation budget: from vacancy thermodynamics to fab spec}
Our defect densities are not free parameters: each Ge$^{2+}$\,$\rightarrow$\,Ge$^{4+}$ oxidation event creates a Ge vacancy (V$_{\mathrm{Ge}}$), the dominant native acceptor in CsGeI$_{3}$ DFT \cite{rMing}, so air exposure sets $N_{t}$ directly. Mapping equilibrium V$_{\mathrm{Ge}}$ concentrations ($N=N_{\mathrm{sites}}\exp(-E^f/kT)$, $N_{\mathrm{sites}}=4.87\times10^{21}$ cm$^{-3}$) onto the consolidated $N_{t}$ ladder (Fig.~25) yields a quantitative fab spec: $E^f\gtrsim0.46$ eV (good glovebox, Ge-rich chemistry) holds the 27.06\% headline; $E^f\approx0.40$ eV (marginal) lands on the realistic 24.40\%; brief air exposure ($E^f\approx0.35$ eV) pushes $p\approx2$[V$_{\mathrm{Ge}}$] past $10^{16}$ cm$^{-3}$, where the $N_{A}$ penalty stays soft ($\le$0.3 pp); ambient processing ($E^f\approx0.30$ eV) falls below the high-oxidation 20.76\%. These are 300~K equilibrium values and hence lower bounds: defects frozen in from typical 100--150~$^{\circ}$C anneals shift every threshold up by 0.11--0.19~eV (holding $10^{14}$ cm$^{-3}$ then needs $E^f\approx0.57$--0.65~eV), and the processing-condition bands are illustrative Ef assumptions, not derived Ge/O chemical-potential regimes. At the Ge-poor limit vacancies reach $\sim$$10^{19}$ cm$^{-3}$ and the absorber goes metallic, as Ming et al.\ predict \cite{rMing} --- consistent with rapid Ge oxidation \cite{rThirumal} and with no reported RbGeI$_{3}$ device to date (phase stability of cubic vs.\ $R3m$ is separately unaddressed here). The mapping is corroborated from three sides: sub-ns TRPL lifetimes invert to $N_{t}\gtrsim10^{17}$, so real Ge films today are 100--1000$\times$ over budget \cite{rGeRev} (S6) --- and a current-state rung at $N_{t}$=$10^{17}$ cm$^{-3}$ on the adopted stack gives 15.53\%, the number a first fabrication should expect to beat; the 1\%-Ge optimum in CsPb(I,Br)$_{3}$ (19.70\% \cite{rGeDope}) marks the same proliferation threshold from the doping side; and passivation chemistries that remove deep levels --- PCCMAI (+3.29 pp \cite{rPCCMAI}), Lewis bases in DFT \cite{rKino} --- are $N_{t}$ reduction in waterfall language. The competing V$_{I}$ deep-trap picture is functional-dependent (benign under HSE+SOC \cite{rZhangPRX}), so this budget rests on V$_{\mathrm{Ge}}$ + Cs$_{I}$ \cite{rZhang24}; CsGeI$_{3}$ energetics transfer to RbGeI$_{3}$ since the A-site carries no band-edge weight. Derivation, bounds and kinetic caveats in S6.

\begin{figure}[htbp]
\centering
\includegraphics[width=\textwidth]{figs/fig39.png}
\caption{Oxidation budget: consolidated PCE vs V$_{\mathrm{Ge}}$ formation energy (dots; equilibrium mapping in text), with processing-condition bands. Dashed extension toward the metallic limit is schematic; TRPL marker from reported sub-ns lifetimes in Ge-containing films (S6).}
\end{figure}


\section{Comparison with Literature}
\subsection{Benchmarking against recent literature}
Table VI compares the optimized FTO/\allowbreak{}TiO$_{2}$/RbGeI$_{3}$/CuI/\allowbreak{}Au device with recently reported RbGeI$_{3}$-based and related lead-free Ge/Sn-based perovskite solar cells. The proposed device achieves a PCE of 27.06\%, within 0.15 pp of the best reported RbGeI$_{3}$ value (27.21\%, Loumachi et al. \cite{r22}), exceeding the 26.47\% of Pathak et al. \cite{r23} (preprint) by 0.59 pp, the 18.60\% of Raj et al. \cite{r39} by 8.46 pp, and the 9.77\% simulation baseline of Pindolia et al. \cite{r14} by 17.29 pp. The performance is competitive with the best tin-based devices (26.38\% CsSnI$_{3}$ \cite{r40}, 25.98\% MASnI$_{3}$ \cite{r41}) while sidestepping Sn$^{2+}$ $\rightarrow$ Sn$^{4+}$ solution oxidation (Ge brings its own Ge$^{2+}$ $\rightarrow$ Ge$^{4+}$ liability, budgeted in Section~VI.H), and exceeds CZTS-HTL devices \cite{r42}.
All comparator values were cross-checked against primary sources (adopting the efficiency implied by reported V$_{\mathrm{OC}}$, J$_{\mathrm{SC}}$, FF where inconsistent); all entries are SCAPS-1D simulations, and the Loumachi string is ITO/\allowbreak{}C$_{60}$/RbGeI$_{3}$/CBTS/\allowbreak{}Ag (Section II).

In short, the optimized device lands within 0.15 pp of the best published RbGeI$_{3}$ simulation (27.21\% \cite{r22}), and stays competitive with the top tin-based lead-free cells, all with an all-inorganic transport stack. With the Cu$_{2}$O-like BSF the ideal point reaches 28.01\% while keeping the cheap TiO$_{2}$/CuI pair (cf.\ Sb$_{2}$S$_{3}$-route 30.35\% \cite{rRahman} --- numerically equal from an independent stack, not the same device --- and IGZO/Cu$_{2}$O 33.84\% \cite{rReza}); the BSF-realistic 21.76\% is the number an experimentalist should target, reported with the full loss decomposition most entries omit. The ML R$^{2}$ (0.907 held out) sits with recent SCAPS+ML reports \cite{rJCE,rMA,rML40}.
 The Pindolia et al. \cite{r14} baseline of 9.77\% leaves a wide gap for improvement. Working RbGeI$_{3}$ solar cells are still rare, so experimental confirmation remains the bottleneck.

\subsection{Gap-plus-alignment control (fixed-$\chi$ protocol)}
To isolate the gap effect, CsGeI$_{3}$ was simulated in the identical stack changing only the bandgap (1.4 $\rightarrow$ 1.6 eV) at fixed $\chi$; a control rerun at 1.4 eV reproduces the RbGeI$_{3}$ point exactly (27.06\%), proving the isolation is computational. The CsGeI$_{3}$ point (Table VII) gives 1.0966 V, 23.1755 mA/cm$^{2}$, 77.3\%, 19.65\%: Cs wins V$_{\mathrm{OC}}$ by 0.035 V through the wider gap while Rb wins J$_{\mathrm{SC}}$ by 7.14 mA/cm$^{2}$ and FF by 6.8 pp. The J$_{\mathrm{SC}}$ gap is dominated by photon harvest --- the 1.6 eV gap excludes the entire 775--886 nm band under this protocol --- while the 0.30 eV cliff's signature is the pinned V$_{\mathrm{OC}}$ ($\sim$1.10 V across a 100$\times$ change in $N_{t}$: 19.65/19.31/17.07\%) and the FF deficit, i.e.\ interface limitation rather than intrinsic robustness. This is therefore evidence about gap-plus-alignment, not A-site chemistry proper --- stated plainly because the same fixed-$\chi$ choice underlies the comparison. The ranking survives derating ($10^{15}$: 24.40 vs 19.31\%; $10^{16}$: 20.76 vs 17.07\%, Fig.~16).


\section{Limitations and Future Work}

\subsection{Limitations of the present study}
The stacked ladder reads 27.06\% ideal $\rightarrow$ 24.28\% (optics) $\rightarrow$ 21.93\% (bulk $10^{15}$) $\rightarrow$ 20.78\% (interfaces $10^{15}$), with $R_{\mathrm{s}}$/$R_{\mathrm{sh}}$ post-corrections (Fig. 16).
Five caveats bound these results. First, the 27.06\% adopted device is an ideal-coupling upper bound at literature-pinned densities; the unconstrained sweep maxima ( Eg 1.3, $N_{\mathrm{C}}$/N$_{\mathrm{V}}$ below $10^{18}$) were rejected by the triple-gate rule, not hidden. Second, the adopted $10^{14}$ densities require well-passivated films: bulk $10^{15}$ ($10^{16}$) drops PCE to 24.40\% (20.76\%), rear interfaces $10^{15}$ ($10^{16}$) to 20.29\% (19.65\%), with J$_{\mathrm{SC}}$ nearly unchanged and V$_{\mathrm{OC}}$/FF collapsing (Fig. 8); holding $10^{14}$ in RbGeI$_{3}$ is hard because Ge$^{2+}$ oxidizes readily to Ge$^{4+}$ in ambient. Third, SCAPS-1D assumes uniform composition and abrupt interfaces that cannot be built exactly. Fourth, the sequential procedure misses cross-interactions by construction; the GA best-found point (28.30\%) is reported as optimizer reach, not a proven maximum, and it sits below the triple-gate density floors so it is not an admissible device under our own rules. Fifth, the oxidation map is an equilibrium bound (300~K values are lower bounds; 100--150~$^{\circ}$C freeze-in shifts thresholds up 0.11--0.19~eV): real oxidation is kinetic, dose-dependent and surface-localized, and CsGeI$_{3}$ energetics transfer to RbGeI$_{3}$ by framework analogy, not calculation.

\subsection{Future work}
The obvious next step is to build and test the optimized device under strict inert processing (<0.1 ppm O2/H2O glovebox), with Ge$^{2+}$ stabilization (reducing additives, trap passivation, alloying, encapsulation) and stability testing. On the modelling side, three extensions are queued: (i) a full 11-parameter global optimization (Bayesian or surrogate-assisted, $\approx$500 evaluations) to confirm the six-gene GA maximum; (ii) a transfer-matrix optical model coupled to the drift-diffusion solution using measured RbGeI$_{3}$ n/k spectra; and (iii) replacing the numerical grading with a realizable composition grade re-optimized in SCAPS.

\section{Conclusion}
We optimized a lead-free RbGeI$_{3}$ perovskite solar cell in SCAPS-1D and converted its central liability --- Ge$^{2+}$ oxidation --- into a quantitative fab spec. From eight ETL--HTL combinations, FTO/\allowbreak{}TiO$_{2}$/RbGeI$_{3}$/CuI/\allowbreak{}Au was selected on durability, cost, and alignment rather than raw efficiency (D2 outscored it 21.10\% to 19.78\% at screening).
After sequential optimization of eleven parameters on one frozen deck, the consolidated device reaches 27.06\% (V$_{\mathrm{OC}}$ = 1.06 V, J$_{\mathrm{SC}}$ = 30.32 mA/cm$^{2}$, FF = 84.11\%), within 0.15 pp of the best prior RbGeI$_{3}$ simulation, with every metric below its Shockley--Queisser limit; the D2/NiO cross-check through the same pipeline reaches 28.16\%, so the reported stack trades the efficiency crown for stability and cost. Joint search returns 28.30\% as optimizer reach below the adoption floors ($E_{\mathrm{g}}$ = 1.36 eV, 62 evaluations), while stacked optical, defect and parasitic losses bound the realistic window to a single number: 20.78\% (21.76\% with BSF). Defect-energy, doping, resistance, mobility, and gap-control sweeps (Sections V.C, V.J--V.L, VII.B) convert the headline into tolerances recomputed on the stacked point: $R_{\mathrm{s}}$ $\le$ 4 (pre-optics), $R_{\mathrm{sh}}$ $\ge$ 10$^{3}$ $\Omega\cdot$cm$^{2}$, $N_{\mathrm{A}}$ $\le$ $10^{16}$ cm$^{-3}$, $\mu$ $\ge$ 15 cm$^{2}$/Vs. Fixing the economical TiO$_{2}$/RbGeI$_{3}$/CuI stack, a p$^{+}$ BSF adds +0.94 pp ideally (27.06 $\rightarrow$ 28.00\%, 28.01\% at the grid optimum) and +0.98 pp on the identical full stack (20.78 $\rightarrow$ 21.76\%) by widening the rear interface gap, with block-level trap attribution and a p$^{+}$ control isolating the offset share; and SHAP sensitivity yields design rules (bulk $N_{t}$, $E_{g}$, $\mu_{n}$) that one-at-a-time sweeps alone cannot provide. Targets for the lab: $N_{A}$(BSF) $\sim$$10^{19}$ cm$^{-3}$ ($\ge$30 nm suffices), bulk $N_{t}$ $\le$$10^{15}$ cm$^{-3}$, and an oxidation budget of $E^f\gtrsim0.46$ eV at 300 K equilibrium (0.57--0.65 eV for anneal freeze-in) for the headline, with 15.53\% marking the rung current films must beat. The pipeline behind these targets reproduces two experimental $\beta$-CsPbI$_{3}$ devices within 2\% in a calibration refit, and its stacked window is placed against a frozen inorganic snapshot as context --- the ideal-to-realistic bridge this field needs before first fabrication.

\section*{Data availability}
All SCAPS decks (initial, H1-adopted, D2, BSF, stacked, attribution controls), run scripts and archived result JSONs listed in Supplementary Section~S7 are available from the authors on request; digitized validation curves and database caches are documented in Supplementary Section~S5.

\begin{thebibliography}{99}
\bibitem{r1} RENEWABLES 2024 Global Status Report (REN21, 2024).
\bibitem{r2} World Energy Outlook 2024 (International Energy Agency, IEA, 2024).
\bibitem{r3} M. A. Green et al., ``Solar cell efficiency tables (Version 64),'' Prog. Photovolt.: Res. Appl. (2024). doi:\url{10.1002/pip.3831.}
\bibitem{r4} R. A. Afre and D. Pugliese, ``Perovskite solar cells: A review of the latest advances in materials, fabrication techniques, and stability enhancement strategies,'' Micromachines 15(2), 192 (2024). doi:\url{10.3390/mi15020192.}
\bibitem{r5} N.-G. Park, M. Grätzel, T. Miyasaka, K. Zhu, and K. Emery, ``Towards stable and commercially available perovskite solar cells,'' Nat. Energy 1(11), 16152 (2016). doi:\url{10.1038/nenergy.2016.152.}
\bibitem{r6} N. J. Jeon, J. H. Noh, Y. C. Kim, W. S. Yang, S. Ryu and S. I. Seok, ``Compositional engineering of perovskite materials for high-performance solar cells,'' Nature 517(7535), 476--480 (2015). doi:\url{10.1038/nature14133.}
\bibitem{r7} I. Ahmed, K. Prakash, and S. M. Mobin, ``Lead-free perovskites for solar cell applications: recent progress, ongoing challenges, and strategic approaches,'' Chem. Commun. 61(37), 6691--6721 (2025). doi:\url{10.1039/D4CC06835A.}
\bibitem{r8} S. Kinza Fatima et al., ``Lead-free perovskites for next-generation applications: a comprehensive computational and data-driven review,'' Mater. Adv. 6(21), 7634--7661 (2025). doi:\url{10.1039/D5MA00681C.}
\bibitem{r9} J. Burschka, N. Pellet, S.-J. Moon, R. Humphry-Baker, P. Gao, M. K. Nazeeruddin, and M. Grätzel, ``Sequential deposition as a route to high-performance perovskite-sensitized solar cells,'' Nature 499(7458), 316--319 (2013). doi:\url{10.1038/nature12340.}
\bibitem{r10} M. M. Lee, J. Teuscher, T. Miyasaka, T. N. Murakami, and H. J. Snaith, ``Efficient hybrid solar cells based on meso-superstructured organometal halide perovskites,'' Science 338(6107), 643--647 (2012). doi:\url{10.1126/science.1228604.}
\bibitem{r11} A. K. Jena, A. Kulkarni, and T. Miyasaka, ``Halide perovskite photovoltaics: Background, status, and future prospects,'' Chem. Rev. 119(5), 3036--3103 (2019). doi:\url{10.1021/acs.chemrev.8b00539.}
\bibitem{r12} P. Afrin, K. Farjana, A. Vumije, and M. N. Uddin, ``An investigation on the impact of temperature variation over the performance of formamidinium tin iodide perovskite solar cell using SCAPS simulation,'' AIP Adv. 14(7) (2024). doi:\url{10.1063/5.0209332.}
\bibitem{r13} S. Samaki, F. Tchangnwa Nya, G. M. Dzifack Kenfack, and A. Laref, ``Materials and interfaces properties optimization for high-efficient and more stable RbGeI$_{3}$ perovskite solar cells: optoelectrical modelling,'' Sci. Rep. 13, 15517 (2023). doi:\url{10.1038/s41598-023-42471-w.}
\bibitem{r14} G. Pindolia, S. M. Shinde, and P. K. Jha, ``Optimization of an inorganic lead-free RbGeI$_{3}$ based perovskite solar cell by SCAPS-1D simulation,'' Sol. Energy 236, 802--821 (2022). doi:\url{10.1016/j.solener.2022.03.053.}
\bibitem{r15} S. Bouazizi, W. Tlili, A. Bouich, B. M. Soucase, and A. Omri, ``Design and efficiency enhancement of FTO/\allowbreak{}PC$_{60}$BM/CsSn0.5Ge0.5I3/\allowbreak{}Spiro-OMeTAD/\allowbreak{}Au perovskite solar cell utilizing SCAPS-1D simulator,'' Mater. Res. Express 9(9), 096402 (2022). doi:\url{10.1088/2053-1591/ac8d52.}
\bibitem{r16} T. Salim, S. Sun, Y. Abe, A. Krishna, A. C. Grimsdale, and Y. M. Lam, ``Perovskite-based solar cells: Impact of materials processing and device architecture,'' J. Mater. Chem. A 3(17), 8943--8969 (2015). doi:\url{10.1039/C4TA05226A.}
\bibitem{r17} P. K. Dash, P. Banarjee, A. Nyayban, and S. Panda, ``Optoelectronic properties and performance optimization for photovoltaic applications of R3m-RbGeX$_{3}$ (X = Cl, Br, I) perovskites: A combined DFT and SCAPS-1D study,'' arXiv:2505.09362 (2025).
\bibitem{r18} N. A. Abdulkareem, ``A DFT study of structural, electronic and optical properties of lead-free and Ge-based cubic perovskite RbGeX$_{3}$ (X = I, Br and Cl),'' Passer J. Basic Appl. Sci. 5(1), 46--51 (2023). doi:\url{10.24271/psr.2023.371069.1187.}
\bibitem{r19} M. Burgelman, P. Nollet, and S. Degrave, ``Modelling polycrystalline semiconductor solar cells,'' Thin Solid Films 361--362, 527--532 (2000). doi:\url{10.1016/S0040-6090(99)00825-1.}
\bibitem{r20} M. Burgelman, M. Debucquoy, K. Decock, J. Verschraegen, and S. Degrave, SCAPS-1D simulation program for thin film solar cells, version 3.3.10 (ELIS, University of Gent, Belgium, 2024).
\bibitem{r22} L. Loumachi, A. Yousfi, O. Saidani, and A. S. Alsubaie, ``Strengthen the power conversion efficiency of solar cell based RbGeI$_{3}$: Numerical approach,'' East Eur. J. Phys. (3), 416--424 (2024). doi:\url{10.26565/2312-4334-2024-3-50}
\bibitem{r23} A. Pathak, K. M. D. Pravalli, R. Madaka, B. Pandey, D. Shankar, C. S. Beera, and D. J. Rao, ``Computational analysis of germanium-based perovskites (MAGeI$_{3}$, FAGeI$_{3}$, CsGeI$_{3}$, RbGeI$_{3}$) for high-efficiency solar cells using SCAPS-1D,'' Research Square (2026). doi:\url{10.21203/rs.3.rs-8749975/}
\bibitem{r24} D. Saikia, J. Bera, A. Betal, and S. Sahu, ``Performance evaluation of an all inorganic CsGeI$_{3}$ based perovskite solar cell by numerical simulation,'' Opt. Mater. 123, 111839 (2022). doi:\url{10.1016/j.optmat.2021.111839.}
\bibitem{r25} R. Pau et al., ``Solution-processed CuI as a hole transport layer for Sn--Pb perovskite solar cells,'' J. Mater. Chem. A 13(48), 42118--42127 (2025). doi:\url{10.1039/D5TA06770G.}
\bibitem{r26} Q. Sun, A. Sadhu, S. Lie, and L. H. Wong, ``Critical review of Cu-based hole transport materials for perovskite solar cells: From theoretical insights to experimental validation,'' Adv. Mater. (2024). doi:\url{10.1002/adma.202402412.}
\bibitem{r27} J. Z. Mbese, ``Advancements in inorganic hole-transport materials for perovskite solar cells: A comparative review,'' Energies 18(9), 2374 (2025). doi:\url{10.3390/en18092374.}
\bibitem{r28} L. T. C. Tuyen, S.-R. Jian, N. T. Tien, and P. H. Le, ``Nanomechanical and material properties of fluorine-doped tin oxide thin films prepared by ultrasonic spray pyrolysis: Effects of F-doping,'' Materials 12(10), 1665 (2019). doi:\url{10.3390/ma12101665.}
\bibitem{r30} A. J. Olasoji and S. H. Im, ``Perspective chapter: TiO$_{2}$ electron transporting layers for perovskite solar cells,'' in Titanium Dioxide -- Uses, Applications, and Advances (IntechOpen, 2024). doi:\url{10.5772/intechopen.1007266.}
\bibitem{r31} K. Valadi, S. Gharibi, R. Taheri-Ledari, S. Akin, A. Maleki, and A. E. Shalan, ``Metal oxide electron transport materials for perovskite solar cells: A review,'' Environ. Chem. Lett. 19(3), 2185--2207 (2021). doi:\url{10.1007/s10311-020-01171-x.}
\bibitem{r32} L. Ji, T. Zhang, and S. Li, ``Band alignment of TiO$_{2}$ by controlling Cl content for high-efficiency perovskite solar cells,'' J. Mater. Chem. C 12(36), 14597--14604 (2024). doi:\url{10.1039/D4TC01495B.}
\bibitem{r33} A. Ghosh et al., ``Solar power conversion: CuI hole transport layer and Ba3NCl3 absorber enable advanced solar cell technology boosting efficiency over 30\%,'' RSC Adv. 14(33), 24066--24081 (2024). doi:\url{10.1039/D4RA03695F.}
\bibitem{r34} T. Minemoto and M. Murata, ``Impact of work function of back contact of perovskite solar cells without hole transport material analyzed by device simulation,'' Curr. Appl. Phys. 14(11), 1428--1433 (2014). doi:\url{10.1016/j.cap.2014.08.002.}
\bibitem{r35} W. Shockley and H. J. Queisser, ``Detailed balance limit of efficiency of p-n junction solar cells,'' J. Appl. Phys. 32(3), 510--519 (1961). doi:\url{10.1063/1.1736034.}
\bibitem{r36} S. Rühle, ``Tabulated values of the Shockley--Queisser limit for single junction solar cells,'' Sol. Energy 130, 139--147 (2016). doi:\url{10.1016/j.solener.2016.02.015.}
\bibitem{r37} C. H. Henry, ``Limiting efficiencies of ideal single and multiple energy gap terrestrial solar cells,'' J. Appl. Phys. 51(8), 4494--4500 (1980). doi:\url{10.1063/1.328272.}
\bibitem{r38} O. J. Sandberg and A. Armin, ``On the effect of surface recombination in thin film solar cells, light emitting diodes and photodetectors,'' arXiv:1905.07291 (2019); published in Synth. Met. 254, 114--121 (2019), doi:\url{10.1016/j.synthmet.2019.06.008.}
\bibitem{r39} M. Raj, A. Aggarwal, A. Kushwaha, and N. Goel, ``Moisture-resistant sustainable solar cell with RbGeI$_{3}$ absorber layer,'' J. Mater. Sci. 60(22), 9176--9196 (2025). doi:\url{10.1007/s10853-025-11002-5.}
\bibitem{r40} H.-J. Park, H. Son, and B.-S. Jeong, ``SCAPS-1D Simulation for Device Optimization to Improve Efficiency in Lead-Free CsSnI$_{3}$ Perovskite Solar Cells,'' Inorganics 12(4), 123 (2024). doi:\url{10.3390/inorganics12040123.}
\bibitem{r41} R. Islam, M. N. Uddin, M. J. Hossen, and A. Begum, ``Photovoltaic performance of lead-free MASnI$_{3}$-based perovskite solar cells with kesterite hole transport materials,'' AIP Adv. 16(5) (2026). doi:\url{10.1063/5.0315071.}
\bibitem{r42} A. C. Piñón Reyes et al., ``Study of a lead-free perovskite solar cell using CZTS as HTL to achieve a 20\% PCE by SCAPS-1D simulation,'' Micromachines 12(12), 1508 (2021). doi:\url{10.3390/mi12121508.}
\bibitem{rReview} V. H. D. Ara\'ujo, A. F. Nogueira, J. C. Trist\~ao and L. J. dos Santos, ``Advances in lead-free perovskite solar cell design via SCAPS-1D simulations,'' RSC Sustain.\ 3, 4314--4335 (2025). doi:\url{10.1039/D5SU00526D.}
\bibitem{rRahman} M. M. Rahman et al., ``Numerical modeling and extensive analysis of an extremely efficient RbGeI$_{3}$-based perovskite solar cell by incorporating a variety of ETL and HTL materials,'' Energy Adv.\ 3(9), 2377--2398 (2024). doi:\url{10.1039/D4YA00323C.}
\bibitem{rReza} M. S. Reza et al., ``Optimizing charge transport layers to enhance the performance of lead-free RbGeI$_{3}$ perovskite solar cells,'' Langmuir 41(11), 7865--7885 (2025).
\bibitem{rJCE} U. F. Ali et al., ``Integrated simulation and machine learning framework for high-performance lead-free RbGeI$_{3}$ perovskite solar cells with WS$_{2}$/CuI transport layers,'' J.\ Comput.\ Electron.\ 25, 37 (2026). doi:\url{10.1007/s10825-025-02476-2.}
\bibitem{rMA} T. M. Khan et al., ``Investigating the impact of non-ideal conditions on the performance of an RbGeI$_{3}$ perovskite solar cell through a combination of SCAPS-1D, machine learning and deep learning approaches,'' Mater.\ Adv.\ (2026). doi:\url{10.1039/D5MA01427A.}
\bibitem{rML40} M. A. Rahman and M. J. Alam, ``Machine learning--guided SCAPS-1D design of lead-free CH$_{3}$NH$_{3}$SnI$_{3}$ perovskite solar cells with 40.17\% simulated PCE,'' J.\ Mater.\ Sci.: Mater.\ Eng.\ (2026). doi:\url{10.1186/s40712-026-00426-9.}
\bibitem{rMing} W. Ming, H. Shi and M.-H. Du, ``Large dielectric constant, high acceptor density, and deep electron traps in perovskite solar cell material CsGeI$_{3}$,'' J.\ Mater.\ Chem.\ A 4, 13852--13858 (2016). doi:\url{10.1039/C6TA04685A.}
\bibitem{rThirumal} K. Thirumal et al., ``Lead-free germanium iodide perovskite materials for photovoltaic applications,'' J.\ Mater.\ Chem.\ A 3, 23829--23832 (2015). doi:\url{10.1039/C5TA05741H.}
\bibitem{rZhang24} C. Zhang et al., ``A systematical study on bands and defects of CsBX$_{3}$ (B = Pb, Sn, Ge; X = Cl, Br, I) perovskite based on first principles,'' Molecules 29(11), 2479 (2024). doi:\url{10.3390/molecules29112479.}
\bibitem{rZhangPRX} J. Zhang, X. Zhang, M. E. Turiansky and C. G. Van de Walle, ``Iodine vacancies do not cause nonradiative recombination in halide perovskites,'' PRX Energy 2, 013008 (2023). doi:\url{10.1103/PRXEnergy.2.013008.}
\bibitem{rPCCMAI} S. Mohammadi, L. J. A. Koster and S. Akbari Nia, ``Enhanced stability and efficiency in Ge--Pb-based perovskite solar cells through PCCMAI-induced Ge defect passivation,'' Energy Environ.\ Mater.\ 8, e70052 (2025). doi:\url{10.1002/eem2.70052.}
\bibitem{rGeDope} R. Chen et al., ``Enhanced stability of Ge-doped CsPb(I$_{1-x}$Br$_{x}$)$_{3}$ perovskite solar cells,'' J.\ Mater.\ Chem.\ A 13, 9822--9829 (2025). doi:\url{10.1039/D4TA08572H.}
\bibitem{rKino} E. Kino et al., ``A theoretical study of the molecular passivation of p-/n-type defects on tin- and germanium-mixed perovskite surfaces using Lewis base/acid,'' Phys.\ Chem.\ Chem.\ Phys.\ 27, 7429--7439 (2025). doi:\url{10.1039/D4CP03817G.}
\bibitem{rGeRev} W. J. Jang, P. J. Park, W.-J. Ong and S. Y. Kim, ``Stable and eco-friendly inorganic lead-free perovskite solar cells: structural, electronic, and defect engineering,'' Commun.\ Mater.\ 7 (2026). doi:\url{10.1038/s43246-026-01158-1.}
\bibitem{rWang} Y. Wang et al., ``Thermodynamically stabilized $\beta$-CsPbI$_{3}$--based perovskite solar cells with efficiencies $>$18\%,'' Science 365, 591--595 (2019). doi:\url{10.1126/science.aav8680.}
\bibitem{rPDB} T. J. Jacobsson et al., The Perovskite Database Project (perovskitedatabase.com, dataset to Feb.\ 2020 + author uploads), via NOMAD (nomad-lab.eu, Perovskite Solar Cells Database, 43,106 entries accessed Oct.\ 2026).
\end{thebibliography}
\end{document}
% TRIMMED-V2
